% from ega3-1-fr.tex
\oldpage[III]{120}
\[
\mathcal F=\varprojlim_k\mathcal F_k
\]
par (I, 10.11.3). D'autre part, il résultera de
(\hyperref[III.3.4.2-fr]{3.4.2}) que chacun des $R^qf_*(\mathcal F)$, étant
cohérent, est naturellement muni d'une structure de
$\mathcal O_{\mathfrak Y}$-Module topologique (I, 10.11.6), et il en est de
même des $R^qf_*(\mathcal F_k)$. Aux homomorphismes canoniques
$\mathcal F\to\mathcal F_k=\mathcal F/\mathcal K^{k+1}\mathcal F$
correspondent canoniquement des homomorphismes
\[
R^qf_*(\mathcal F)\longrightarrow R^qf_*(\mathcal F_k)
\]
qui sont nécessairement continus pour les structures de
$\mathcal O_{\mathfrak Y}$-Module topologique précédentes (I, 10.11.6), et
forment un système projectif, donnant à la limite un homomorphisme canonique
fonctoriel
\[
R^qf_*(\mathcal F)\longrightarrow\varprojlim_kR^qf_*(\mathcal F_k)
\tag{3.4.2.2}\label{III.3.4.2.2-fr}
\]
qui sera un homomorphisme continu de $\mathcal O_{\mathfrak Y}$-Modules
topologiques. Nous allons démontrer en même temps que
(\hyperref[III.3.4.2-fr]{3.4.2}) le

\begin{corollary}[3.4.3]
\label{III.3.4.3-fr}
Chacun des homomorphismes (\hyperref[III.3.4.2.2-fr]{3.4.2.2}) est un
isomorphisme topologique. En outre, si $\mathfrak Y$ est noethérien, le
système projectif
$(R^qf_*(\mathcal F/\mathcal K^{k+1}\mathcal F))_{k\geq0}$ satisfait à la
condition (ML) (0, 13.1.1).
\end{corollary}

Nous commencerons par établir (\hyperref[III.3.4.2-fr]{3.4.2}) et
(\hyperref[III.3.4.3-fr]{3.4.3}) lorsque $Y$ est un schéma affine
formel noethérien (I, 10.4.1)~:

\begin{corollary}[3.4.4]
\label{III.3.4.4-fr}
Sous les hypothèses de (\hyperref[III.3.4.2-fr]{3.4.2}), supposons en outre
que $\mathfrak Y=\operatorname{Spf}(A)$, où $A$ est un anneau adique
noethérien. Soit $\mathfrak J$ un idéal de définition de $A$, et posons
$\mathcal F_k=\mathcal F/\mathfrak J^{k+1}\mathcal F$ pour $k\geq0$. Alors
les $H^n(\mathfrak X,\mathcal F)$ sont des $A$-modules de type fini~; le
système projectif $(H^n(\mathfrak X,\mathcal F_k))_{k\geq0}$ satisfait à la
condition (ML) pour tout $n$~; si on pose
\[
N_{n,k}=\operatorname{Ker}\bigl(H^n(\mathfrak X,\mathcal F)
\longrightarrow H^n(\mathfrak X,\mathcal F_k)\bigr)
\tag{3.4.4.1}\label{III.3.4.4.1-fr}
\]
(égal aussi à
$\operatorname{Im}(H^n(\mathfrak X,\mathfrak J^{k+1}\mathcal F)
\to H^n(\mathfrak X,\mathcal F))$ par la suite exacte de cohomologie), les
$N_{k,n}$ définissent sur $H^n(\mathfrak X,\mathcal F)$ une filtration
$\mathfrak J$-bonne (0, 13.7.7)~; enfin, l'homomorphisme canonique
\[
H^n(\mathfrak X,\mathcal F)
\longrightarrow\varprojlim_kH^n(\mathfrak X,\mathcal F_k)
\tag{3.4.4.2}\label{III.3.4.4.2-fr}
\]
est un isomorphisme topologique pour tout $n$ (le premier membre étant muni
de la topologie $\mathfrak J$-adique, les $H^n(\mathfrak X,\mathcal F_k)$ de
la topologie discrète).
\end{corollary}

Posons
\[
S=\operatorname{gr}(A)=\bigoplus_{k\geq0}
\mathfrak J^k/\mathfrak J^{k+1},\qquad
\mathcal M=\operatorname{gr}(\mathcal F)=\bigoplus_{k\geq0}
\mathfrak J^k\mathcal F/\mathfrak J^{k+1}\mathcal F.
\tag{3.4.4.3}\label{III.3.4.4.3-fr}
\]
On sait que $\mathfrak J^\Delta$ est un Idéal de définition de $\mathfrak Y$
(I, 10.3.1)~; soient
$\mathcal K=f^*(\mathfrak J^\Delta)\mathcal O_{\mathfrak X}$,
$X_0=(\mathfrak X,\mathcal O_{\mathfrak X}/\mathcal K)$,
$Y_0=(\mathfrak Y,\mathcal O_{\mathfrak Y}/\mathfrak J^\Delta)
=\operatorname{Spec}(A_0)$ avec $A_0=A/\mathfrak J$. Il est clair que les
$\mathcal M_k=\mathfrak J^k\mathcal F/\mathfrak J^{k+1}\mathcal F$ sont des
$\mathcal O_{X_0}$-Modules cohérents (I, 10.11.3). Considérons d'autre part la
$\mathcal O_{X_0}$-Algèbre graduée quasi-cohérente
\[
\mathcal S=\mathcal O_{X_0}\otimes_{A_0}S
=\operatorname{gr}(\mathcal O_{\mathfrak X})
=\bigoplus_{k\geq0}\mathcal K^k/\mathcal K^{k+1}.
\tag{3.4.4.4}\label{III.3.4.4.4-fr}
\]

L'hypothèse que $\mathcal F$ est un $\mathcal O_{\mathfrak X}$-Module de type
fini entraîne d'abord que $\mathcal M$ est

